\documentclass[a4paper,11pt,oneside]{book}
% preamble_build.tex: copy of Maximus's preamble.tex for the DE v3 build.
% The blocks marked [BUILD COPY CHANGE] are the earlier build fixes.
% The blocks marked [V3 ADDITION] are new in this version (example box,
% spacing). Everything else is unchanged from preamble.tex.

% For margins
\usepackage[a4paper, margin = 20mm]{geometry}

\usepackage[Lenny]{fncychap}
% Rejne, Glenn, Lenny

\usepackage{etoc}

% [V3 ADDITION] cancel: used to show the dx's "cancelling" in Section 11.3
\usepackage{cancel}
\renewcommand{\CancelColor}{\color{red}}

% For enumeration & math symbols
\usepackage{enumerate}
\usepackage{amsmath,amsfonts,amssymb,amsthm,graphicx,lipsum}

% For diagrams, plots and boxes
\usepackage{pgfplots}
\pgfplotsset{compat=1.18}
\usepackage{tikz}
\usepackage[most]{tcolorbox}
\tcbuselibrary{skins}
\usepackage{xcolor}
\pgfplotsset{
    standard/.style={clip = false,
    axis line style = thick,
    trig format=rad,
    axis lines = middle,
    every axis x label/.style={at={(current axis.right of origin)},anchor= west},
    every axis y label/.style={at={(current axis.above origin)},anchor=south}
    }
}

% For GC font
\usepackage{fontspec}
% [BUILD COPY CHANGE 1] The TI-84 key font is loaded only if the file exists
% (checked in ./Fonts/... and ../Fonts/...). Otherwise \ti{} falls back to a
% small boxed sans-serif key. Known letters of the TI font are mapped to their
% key names in the fallback: o = [y=], z = [mode], f = [L3].
\newcommand{\tikey}[1]{\tcbox[on line, enhanced, boxsep=0pt, left=2pt, right=2pt,
  top=1pt, bottom=1pt, boxrule=0.5pt, arc=1.5pt, colback=white,
  colframe=black!70]{\sffamily\footnotesize\strut #1}}
\ExplSyntaxOn
\NewDocumentCommand{\tifallback}{m}{
  \str_case:nnF {#1}
    { {o}{\tikey{y=}} {z}{\tikey{mode}} {f}{\tikey{L3}} }
    { \tikey{#1} }
}
\ExplSyntaxOff
\def\tifontcase{0}
\IfFileExists{Fonts/TI-84 Plus CE Key Fonts/TI84PlusCEKeys.ttf}
  {\def\tifontcase{1}}{}
\IfFileExists{../Fonts/TI-84 Plus CE Key Fonts/TI84PlusCEKeys.ttf}
  {\def\tifontcase{2}}{}
\ifcase\tifontcase\relax
  % 0: font not found, use boxed fallback keys
  \newcommand{\ti}[1]{\tifallback{#1}}
\or
  % 1: original location (as in preamble.tex)
  \newfontfamily\tifont[
    Path = {Fonts/TI-84 Plus CE Key Fonts/},
    UprightFont = {TI84PlusCEKeys.ttf}
  ]{}
  \newcommand{\ti}[1]{{\tifont #1}}
\or
  % 2: one level up (when compiling inside draft-ch1/)
  \newfontfamily\tifont[
    Path = {../Fonts/TI-84 Plus CE Key Fonts/},
    UprightFont = {TI84PlusCEKeys.ttf}
  ]{}
  \newcommand{\ti}[1]{{\tifont #1}}
\fi

% [FONT SETUP copied from h2-math-notes/aoi/notes/preamble_build.tex (DI notes)]
% [AoI items 17/18/37] GC DISPLAY FONT. Every GC key, command, menu item and
% screen line goes through \gcfont, so the font can be swapped HERE in one place.
%   \gc{fnInt(}      command / menu item / stored variable in the display font
%   \gckey{2nd}      a calculator key (display font in a small key outline)
%   \gcscreen{...}   a calculator screen (display font on a grey-green LCD panel)
% The font file lives in ../Fonts (shared by notes, slides and worksheet).
\def\gcfontfile{VT323-Regular.ttf}
\IfFileExists{Fonts/\gcfontfile}{%
  \newfontfamily\gcfont[Path=Fonts/, Scale=1.32]{\gcfontfile}}{%
  \newcommand\gcfont{\ttfamily}}
\newcommand{\gc}[1]{\mbox{\gcfont #1}}
\newcommand{\gckey}[1]{\tcbox[on line, enhanced, boxsep=0pt, left=2pt, right=2pt,
  top=0.6pt, bottom=0.2pt, boxrule=0.5pt, arc=1.5pt, colback=black!4,
  colframe=black!65]{\gcfont\strut #1}}
\definecolor{lcd}{RGB}{214,222,204}
\newcommand{\gcscreen}[1]{\tcbox[on line, enhanced, boxsep=1pt, left=4pt, right=4pt,
  top=2pt, bottom=2pt, boxrule=1.2pt, arc=1pt, colback=lcd, colframe=black!70]{%
  \gcfont\begin{tabular}{@{}l@{}}#1\end{tabular}}}
% integral template as shown on the TI-84 Plus CE (MathPrint): \gcint{0}{2}{e^(-X²)}{X}
\newcommand{\gcint}[4]{{\gcfont\raisebox{-0.35ex}{\Large ∫}\hspace{-0.15em}%
  \raisebox{-0.2ex}{\begin{tabular}{@{}l@{}}\footnotesize #2\\[-1.3ex]\footnotesize #1\end{tabular}}%
  \hspace{0.1em}(#3)d#4}}
% old key macro kept as an alias so nothing breaks
\renewcommand{\ti}[1]{\gckey{#1}}

% SI units
\usepackage{siunitx}

% setspace
%\usepackage{setspace}
%\onehalfspacing

% For hyperlinks within and out of document
\usepackage{hyperref}
\hypersetup{
    colorlinks = true,
    linkcolor = black,
    urlcolor = blue,
}   

% For captions for pics
\usepackage{capt-of}


% tcolorbox custom environments
% Definition ========================================================
\newtcolorbox[auto counter,number within=section]{definition}[2][drop shadow]{%
enhanced, before skip=5mm, after skip=6mm, % [V3 ADDITION] spacing
colback=blue!5!white,colframe=blue!50!black,fonttitle=\bfseries,
title= Definition~\thetcbcounter. #2,#1}

% Theorem ==========================================================
\newtcolorbox[auto counter,number within=section]{theorem}[2][drop shadow]{%
enhanced, before skip=5mm, after skip=6mm, % [V3 ADDITION] spacing
colback=black!5!white,colframe=black,fonttitle=\bfseries,
title= Theorem~\thetcbcounter. #2,#1}

% Caution box ======================================================
\newtcolorbox{caution}[1][]{enhanced,
before skip=5mm,after skip=6mm, % [V3 ADDITION] was 2mm/3mm
boxrule=1pt,left=11mm,right=2mm,top=3mm,bottom=3mm,
colback=yellow!50,
colframe=yellow!20!black,
sharp corners,rounded corners=southeast,arc is angular,arc=3mm,
underlay={%
\path[fill=tcbcolback!80!black] ([yshift=3mm]interior.south east)--++(-0.4,-0.1)--++(0.1,-0.2);
\path[draw=tcbcolframe,shorten <=-0.05mm,shorten >=-0.05mm] ([yshift=3mm]interior.south east)--++(-0.4,-0.1)--++(0.1,-0.2);
\path[fill=yellow!50!black,draw=none] (interior.south west) rectangle node[white]{\Huge\bfseries !} ([xshift=9mm]interior.north west);
},
drop fuzzy shadow,#1}

% Question box ======================================================
\newtcolorbox{question}[1][]{enhanced,
before skip=5mm,after skip=6mm, % [V3 ADDITION] was 2mm/3mm
boxrule=1pt,left=11mm,right=2mm,top=3mm,bottom=3mm,
colback=blue!20!white,
colframe=blue!20!black,
sharp corners,rounded corners=southeast,arc is angular,arc=3mm,
underlay={%
\path[fill=tcbcolback!80!black] ([yshift=3mm]interior.south east)--++(-0.4,-0.1)--++(0.1,-0.2);
\path[draw=tcbcolframe,shorten <=-0.05mm,shorten >=-0.05mm] ([yshift=3mm]interior.south east)--++(-0.4,-0.1)--++(0.1,-0.2);
\path[fill=blue!50!black,draw=none] (interior.south west) rectangle node[white]{\Huge\bfseries ?} ([xshift=9mm]interior.north west);
},
drop fuzzy shadow,#1}

% Information box ======================================================
\newtcolorbox{info}[1][]{enhanced,
before skip=5mm,after skip=6mm, % [V3 ADDITION] was 2mm/3mm
boxrule=1pt,left=11mm,right=2mm,top=3mm,bottom=3mm,
colback=cyan!20!white,
colframe=cyan!20!black,
sharp corners,rounded corners=southeast,arc is angular,arc=3mm,
underlay={%
\path[fill=tcbcolback!80!black] ([yshift=3mm]interior.south east)--++(-0.4,-0.1)--++(0.1,-0.2);
\path[draw=tcbcolframe,shorten <=-0.05mm,shorten >=-0.05mm] ([yshift=3mm]interior.south east)--++(-0.4,-0.1)--++(0.1,-0.2);
\path[fill=cyan!50!black,draw=none] (interior.south west) rectangle node[white]{\Huge\bfseries i} ([xshift=9mm]interior.north west);
},
drop fuzzy shadow,#1}

% Mistake box ======================================================
\newtcolorbox{mistake}[1][]{enhanced,
before skip=5mm,after skip=6mm, % [V3 ADDITION] was 2mm/3mm
boxrule=1pt,left=11mm,right=2mm,top=3mm,bottom=3mm,
colback=red!20!white,
colframe=red!20!black,
sharp corners,rounded corners=southeast,arc is angular,arc=3mm,
underlay={%
\path[fill=tcbcolback!80!black] ([yshift=3mm]interior.south east)--++(-0.4,-0.1)--++(0.1,-0.2);
\path[draw=tcbcolframe,shorten <=-0.05mm,shorten >=-0.05mm] ([yshift=3mm]interior.south east)--++(-0.4,-0.1)--++(0.1,-0.2);
\path[fill=red!50!black,draw=none] (interior.south west) rectangle node[white]{\Huge\bfseries $\times$} ([xshift=9mm]interior.north west);
},
drop fuzzy shadow,#1}



% [V3 ADDITION] Worked example box ==================================
% Own counter, numbered within the section (Example 11.2.1, 11.2.2, ...).
% For per-subsection numbering (11.2.1.1) change `section' to `subsection'. Blue `question' boxes are kept for thought-provoking questions
% only; worked examples always use this box. Usage:
%   \begin{example} ...question text... \end{example}
%   \solution{...}
\newtcolorbox[auto counter,number within=section]{example}[1][]{%
enhanced, breakable=false, before skip=5mm, after skip=4mm,
colback=teal!5!white, colframe=teal!55!black, boxrule=0.8pt,
left=3mm, right=3mm, top=2mm, bottom=2mm,
sharp corners, rounded corners=northwest, arc=3mm,
fonttitle=\bfseries, coltitle=white,
title=Example~\thetcbcounter, attach boxed title to top left={yshift=-2mm, xshift=3mm},
boxed title style={colback=teal!55!black, sharp corners, rounded corners=northwest, arc=2mm},
drop fuzzy shadow, #1}

% [V3 ADDITION] Spacing between sections, subsections, parts, paragraphs =
% (fncychap's chapter heading is untouched; only \section and below.)
\makeatletter
\renewcommand\section{\@startsection{section}{1}{\z@}%
  {-2ex \@plus -.5ex \@minus -.2ex}{3.2ex \@plus .3ex}%
  {\normalfont\LARGE\bfseries\color{blue!50!black}}}
\renewcommand\subsection{\@startsection{subsection}{2}{\z@}%
  {-1.5ex \@plus -.5ex \@minus -.2ex}{2.6ex \@plus .3ex}%
  {\normalfont\Large\bfseries}}
\renewcommand\subsubsection{\@startsection{subsubsection}{3}{\z@}%
  {-4ex \@plus -1ex \@minus -.5ex}{1.6ex \@plus .2ex}%
  {\normalfont\large\bfseries\color{teal!50!black}}}
\makeatother
\setlength{\parindent}{0pt}
\setlength{\parskip}{0.9ex plus 0.3ex minus 0.2ex}

% Headers and Footers
\usepackage{fancyhdr,lastpage}
% [V3 R2] grey left-margin line on the out-of-syllabus pages (item 13)
\usepackage{eso-pic}
\newcommand{\oosmargin}{\AddToShipoutPictureBG{\put(12mm,0){\color{gray!75}\rule{1.6pt}{\paperheight}}}}
\pagestyle{fancy}
\lhead{Maximus Pung Jun}
% [BUILD COPY CHANGE 2] Header bug fix: \chaptername is the word "Chapter",
% so the original header printed "Chapter 1: Chapter". \chaptermark now stores
% the bare chapter title in \leftmark, and the header uses \leftmark.
\renewcommand{\chaptermark}[1]{\markboth{#1}{}}
\rhead{Chapter \thechapter: \leftmark}
% \cfoot{Page \thepage\  of \pageref{LastPage}} % buggy, does not appear on chapter pages.

% [V3 R2] tick and cross markers for mistake boxes (G3)
\newcommand{\wrong}{\textcolor{red!70!black}{\ensuremath{\boldsymbol{\times}}}}
\newcommand{\cmark}{\textcolor{green!50!black}{\ensuremath{\checkmark}}}

\newcommand{\xmark}{\wrong}   % alias: red cross
% [V3 R2] out-of-syllabus helpers (DE notes item 13)
%   \oosmargin    : grey line in the left margin on every following page
%   \oosmarginoff : stop the margin line
%   \oosnote      : grey banner "Out of syllabus. For your interest only."
%   outofsyllabus : environment = new page + margin line + banner; ends with
%                   \clearpage and switches the margin line off again
\newcommand{\oosmarginoff}{\ClearShipoutPictureBG}
\newcommand{\oosnote}[1][Nothing in this section is examined in H2 Mathematics.]{\par\noindent\colorbox{gray!15}{\parbox{\dimexpr\textwidth-2\fboxsep}{\textbf{Out of syllabus. For your interest only.} #1}}\par}
\newenvironment{outofsyllabus}{\clearpage\oosmargin}{\clearpage\oosmarginoff}

% Notation as in Maximus's notes: italic d, display style everywhere (G1).
% Fractions inside exponents stay small, so never use these macros inside ^{}.
\newcommand{\dydx}{\dfrac{dy}{dx}}
\newcommand{\dvdx}{\dfrac{dv}{dx}}
\newcommand{\dudx}{\dfrac{du}{dx}}
\newcommand{\ddydx}{\dfrac{d^{2}y}{dx^{2}}}
\newcommand{\ddt}[1]{\dfrac{d#1}{dt}}
% Small syllabus flag (no longer used in the DE notes, kept for other chapters)
\newcommand{\syllabusflag}[1]{\par\noindent{\footnotesize\color{red!70!black}\textbf{Syllabus note.} #1}\par}


% Teacher's copy toggle ================================================
% [BUILD COPY CHANGE 3] \ifx compares the whole word (the original
% \if\conform test was always false). DE_notes_student.tex predefines
% \conform as HideSolutions.
% A student copy defines \conform as HideSolutions BEFORE loading this preamble.
\providecommand\conform{ShowSolutions}
\newcommand\ShowSolutionsKey{ShowSolutions}
% [V3 ADDITION] In the student copy the hidden solution leaves stretchable
% blank space (\vfill), so the free space on each page is shared out as
% working space below the examples instead of collecting at the bottom.
\newcommand\solution[1]{\ifx\conform\ShowSolutionsKey
  \par\noindent\textbf{Solution.} #1\par\medskip\else\par\vspace{\workspace}\fi}
% [V3 R2] minimum working space left under each example in the student copy
\newlength\workspace\setlength\workspace{3cm plus 1fill}


\newcommand{\NN}{\mathbb{N}}
\newcommand{\ZZ}{\mathbb{Z}}
\newcommand{\QQ}{\mathbb{Q}}
\newcommand{\RR}{\mathbb{R}}
\newcommand{\CC}{\mathbb{C}}


% =====================================================================
\usepackage{pifont}
% preamble_local.tex: extra packages and styles for the Vectors chapter
% ONLY. Loaded after preamble_build.tex (a verbatim copy of the shared
% book preamble from the PnC / DE build, so that it can be re-synced by
% copying it over again).
% =====================================================================
\usepackage{array, booktabs}
\usetikzlibrary{arrows.meta, calc, positioning, 3d, decorations.pathreplacing, angles, quotes}
% =====================================================================
% vec_common.tex: shared by the notes (via notes/preamble_local.tex) and
% the slides (via slides/slides_local.tex) for the Vectors chapter.
%
%   CHAPTER NUMBER: change the 13 below and rebuild. The notes set the
%   chapter counter from it, and every slide frame title, Example/Quiz
%   number and box number on the slides uses \chapnum.
% =====================================================================
\def\chapnum{13}

% ---- answers -------------------------------------------------------------
% \ans{key}{value} just prints the value. checks/verify.py reads every
% \ans in the notes and slides and checks it with sympy.
\newcommand{\ans}[2]{#2}

% Column vector, the form used in MF27 and in the examples.
% pmatrix cells are text style, so \dfrac alone still looks small and the
% rows collide. Force display style in every cell and open the rows up.
% Each cell gets a strut as tall as its own content plus \colpad above and
% below, so rows of \dfrac never touch, while integer rows stay compact.
\newlength{\colpad}\setlength{\colpad}{0.55ex}
\newcommand{\colcell}[1]{%
  \setbox0=\hbox{$\displaystyle #1$}%
  \rule[-\dimexpr\dp0+\colpad\relax]{0pt}{\dimexpr\ht0+\dp0+2\colpad\relax}%
  \box0}
\newcommand{\col}[3]{%
  \begingroup
  \renewcommand{\arraystretch}{1.35}%
  \begin{pmatrix}
    \colcell{#1}\\
    \colcell{#2}\\
    \colcell{#3}
  \end{pmatrix}%
  \endgroup
}

% \place{Name}{x}{y}{z} defines a TikZ coordinate. The parentheses are
% inside the macro, so TikZ sees them only after expansion, and
% checks/verify.py can read the drawn coordinates back out.
\newcommand{\place}[4]{\coordinate (#1) at (#2,#3,#4);}

% ---- 3D pictures ----------------------------------------------------------
% Oblique axes: x to the right and slightly down, y out of the page,
% z straight up. Use scale to fit the frame.
\tikzset{
  threeD/.style={
    x={(0.78cm,-0.20cm)},
    y={(0.36cm,0.26cm)},
    z={(0cm,0.78cm)},
    line cap=round, line join=round,
    >=Stealth
  },
  varr/.style={->, thick},
  pt/.style={circle, fill=black, inner sep=1.15pt},
  plane/.style={fill=blue!12, draw=blue!45!black, thick, fill opacity=0.85},
  hid/.style={dashed, black!55}
}
% vecfig.tex: shared 3D diagram style (tikz-3dplot) for Vectors I / II notes and slides.
\usepackage{tikz-3dplot}
\tdplotsetmaincoords{62}{55}
\definecolor{vA}{RGB}{31,90,180}   % first vector / line 1
\definecolor{vB}{RGB}{210,90,20}   % second vector / line 2
\definecolor{vC}{RGB}{0,128,110}   % result / third
\definecolor{vN}{RGB}{140,40,150}  % normals
\definecolor{vP}{RGB}{200,215,240} % plane 1
\definecolor{vQ}{RGB}{245,215,190} % plane 2
\definecolor{vR}{RGB}{200,235,215} % plane 3
\tikzset{
  v3/.style={tdplot_main_coords, line cap=round, line join=round, >=Stealth, scale=#1},
  v3/.default=1,
  ax/.style={->, black!50, thin},
  va/.style={->, very thick, vA},
  vb/.style={->, very thick, vB},
  vc/.style={->, very thick, vC},
  vn/.style={->, very thick, vN},
  ln/.style={very thick},
  dsh/.style={dashed, black!55},
  pp/.style={fill=vP, fill opacity=0.55, draw=vA!70!black, thin},
  pq/.style={fill=vQ, fill opacity=0.55, draw=vB!70!black, thin},
  pr/.style={fill=vR, fill opacity=0.55, draw=vC!70!black, thin},
  dot/.style={circle, fill=black, inner sep=1.2pt},
}
% \vaxes{xlen}{ylen}{zlen}
\newcommand{\vaxes}[3]{%
  \draw[ax] (0,0,0) -- (#1,0,0) node[anchor=west]{$x$};
  \draw[ax] (0,0,0) -- (0,#2,0) node[anchor=south west]{$y$};
  \draw[ax] (0,0,0) -- (0,0,#3) node[anchor=south]{$z$};
  \node[anchor=east] at (0,0,0) {\scriptsize $O$};}
% \rightmark{corner}{along1}{along2} small right-angle square
\newcommand{\rtang}[3]{\draw[thin] ($(#1)!0.25cm!(#2)$) -- ($($(#1)!0.25cm!(#2)$)+($(#1)!0.25cm!(#3)$)-(#1)$) -- ($(#1)!0.25cm!(#3)$);}
\newenvironment{vis}{\par\medskip\begin{center}}{\end{center}\medskip}
% \vang{vertex}{p1}{p2}{label}: small angle arc drawn on the page
\newcommand{\vang}[4]{%
  \coordinate (vg1) at ($(#1)!0.55cm!(#2)$); \coordinate (vg2) at ($(#1)!0.55cm!(#3)$);
  \coordinate (vgm) at ($(vg1)!0.5!(vg2)$);
  \draw[vC, thick] (vg1) .. controls ($(#1)!0.72cm!(vgm)$) .. (vg2);
  \node[vC] at ($(#1)!0.95cm!(vgm)$) {#4};}


% A little air between the parts of an example, and between lines of a list.
\setlength{\itemsep}{0.85ex}
\setlength{\parsep}{0.35ex}

% Student copy: minimum working space under each example. 2.2cm (instead of
% the book's 3cm) is meant to keep the student copy on the same page breaks
% as the teacher copy.
\setlength\workspace{2.2cm plus 1fill}

% =====================================================================
% prelim_fonts.tex: fonts of the real 9758 prelim papers (Word + MathType; pdffonts on
% ACJC 2026 P1/P2, ACJC 2025 P1, JPJC 2026, RVHS 2026, CJC/DHS/TJC/NYJC/HCI/EJC 2025):
%   text Times New Roman; cover/footer Arial; RI page footer Calibri (here Carlito);
%   maths = MathType: Times New Roman letters/digits (italic variables; upright digits,
%   function names, d, e, f), SymbolMT for the integral, sigma, pi and = + < >,
%   scripts 7 pt on 12 pt (about 58 %), tight operator spacing.
% Candidate C (teacher-approved 13 Oct 2026): the real Times New Roman and Arial
% (ttf-mscorefonts-installer) + URW Standard Symbols PS (fonts-urw-base35, a SymbolMT
% clone with the same Symbol encoding) for the integral, sigma, pi, = + < >;
% TeX Gyre Termes Math only for what MathType draws itself (radicals, stretchy brackets,
% fraction bars, arrows). House style: upright d (\mathrm{d}x), upright e and f.
% Loaded once by di_style / de_core (HW sheets) and prelim_style (Tests, solutions).
% =====================================================================
\ifdefined\prelimfontsloaded\endinput\fi
\def\prelimfontsloaded{1}
\edef\pfrestoreat{\catcode`\noexpand\@=\the\catcode`\@\relax}\makeatletter
% MathType script sizes: 7 pt superscripts on 12 pt text (58 %), 5 pt second level
\DeclareMathSizes{12}{12}{7}{5}
\DeclareMathSizes{11}{11}{6.5}{5}
\DeclareMathSizes{10.95}{10.95}{6.5}{5}
\DeclareMathSizes{10}{10}{6}{5}
\DeclareMathSizes{9}{9}{5.5}{5}
\DeclareMathSizes{14.4}{14.4}{8.5}{6}
\setmainfont{Times New Roman}
\setsansfont{Arial}
\newfontfamily\coverfont{Arial}
\newfontfamily\calibrifont{Carlito}
\usepackage{unicode-math}
\setmathfont{TeX Gyre Termes Math}
\setmathfont[range={up/{latin,Latin,num}}]{Times New Roman}
\setmathfont[range={it/{latin,Latin}}]{Times New Roman Italic}
\setmathfont[range={bfup/{latin,Latin,num}}]{Times New Roman Bold}
\setmathfont[range={"002B,"003D,"003C,"003E}]{Standard Symbols PS}
\newfontfamily\pf@sym{Standard Symbols PS}
\newfontfamily\pf@symi[FakeSlant=0.21]{Standard Symbols PS}
% MathType-like spacing (tighter than TeX around = + -)
\AtBeginDocument{\thickmuskip=2.5mu plus 1mu\medmuskip=1.5mu plus 1mu\thinmuskip=3mu
% Symbol integral and sigma (18 pt at 12 pt text, as MathType), slanted Symbol pi
 \renewcommand{\int}{\mathop{\vcenter{\hbox{\pf@sym\fontsize{\fpeval{1.5*\f@size}}{0pt}\selectfont\char"F2}}}\nolimits}%
 \renewcommand{\sum}{\mathop{\vcenter{\hbox{\pf@sym\fontsize{\fpeval{1.5*\f@size}}{0pt}\selectfont\char"E5}}}\limits}%
 \renewcommand{\pi}{\mathord{\text{\pf@symi\char"70}}}}
\pfrestoreat

\usepackage{needspace,longtable,multicol,enumitem,array}
\providecommand{\vv}[1]{\overrightarrow{#1}}\providecommand{\cvv}[3]{\begin{pmatrix}#1\\#2\\#3\end{pmatrix}}
\raggedbottom
\setlength{\parindent}{0pt}

\fancyhf{}\lhead{\small Vectors I Test}\rhead{\small Compiled by Maximus Pung Jun}
\cfoot{\small\thepage}
\renewcommand{\headrulewidth}{0.4pt}
\newcounter{wq}
\newcommand{\Q}[1]{\par\needspace{5\baselineskip}\medskip\stepcounter{wq}\noindent\textbf{\thewq}\quad #1\par}
\begin{document}\thispagestyle{fancy}
\setlength{\abovedisplayskip}{4pt}\setlength{\belowdisplayskip}{4pt}\setlength{\abovedisplayshortskip}{2pt}\setlength{\belowdisplayshortskip}{3pt}\setlength{\jot}{2pt}
{\LARGE\scshape Vectors I Test --- Solutions}\par\smallskip
{Compiled by Maximus Pung Jun}\par\smallskip
{\small Full sourced questions. Give exact answers unless otherwise stated.}\par\medskip
\par\noindent\begin{minipage}{\linewidth}
\Q{\textcolor{black!60}{\small[DHS 2023 Promo Q5]}\par\smallskip
With reference to the origin $O$, the points $A$ and $B$ have coordinates $(1,2,0)$ and $(-5,1,0)$.
\begin{enumerate}[label=\textbf{(\alph*)},leftmargin=2.2em,itemsep=3pt]
\item Find the angle $AOB$.\hfill[3]
\end{enumerate}
The vectors $\vv{OA}$ and $\vv{OB}$ are projected along the direction of $\cvv{1}{2}{2}$ and $\cvv{2}{1}{2}$ to form the vectors $\vv{O'A'}$ and $\vv{O'B'}$ respectively.
\begin{enumerate}[label=\textbf{(\alph*)},leftmargin=2.2em,itemsep=3pt,start=2]
\item Find the vectors $\vv{O'A'}$ and $\vv{O'B'}$.\hfill[3]
\item Find the exact area of the triangle $O'A'B'$.\hfill[2]
\end{enumerate}}

\end{minipage}\par
\par\medskip\hrule\medskip
\textbf{Solution}\par\smallskip
\textbf{(a)} $\vv{OA}=\cvv{1}{2}{0}$ and $\vv{OB}=\cvv{-5}{1}{0}$. Let $\theta=\angle AOB$.
\begin{align*}
\vv{OA}\cdot\vv{OB}&=|\vv{OA}||\vv{OB}|\cos\theta\\
\cos\theta&=\frac{\cvv{1}{2}{0}\cdot\cvv{-5}{1}{0}}{\sqrt{1^2+2^2}\sqrt{5^2+1^2}}
=\frac{-5+2}{\sqrt5\sqrt{26}}=-\frac3{\sqrt5\sqrt{26}}\\
\theta&=105.3^\circ.
\end{align*}
\textbf{(b)}
\begin{align*}
\vv{O'A'}&=\frac{\vv{OA}\cdot\cvv{1}{2}{2}}{\left|\cvv{1}{2}{2}\right|}\frac{\cvv{1}{2}{2}}{\left|\cvv{1}{2}{2}\right|}
=\frac{\cvv{1}{2}{0}\cdot\cvv{1}{2}{2}}{(\sqrt{1^2+2^2+2^2})^2}\cvv{1}{2}{2}=\frac59\cvv{1}{2}{2}\\
\vv{O'B'}&=\frac{\vv{OB}\cdot\cvv{2}{1}{2}}{\left|\cvv{2}{1}{2}\right|}\frac{\cvv{2}{1}{2}}{\left|\cvv{2}{1}{2}\right|}
=\frac{\cvv{-5}{1}{0}\cdot\cvv{2}{1}{2}}{(\sqrt{2^2+1^2+2^2})^2}\cvv{2}{1}{2}=-\cvv{2}{1}{2}.
\end{align*}
\textbf{(c)}
\begin{align*}
\text{Area of triangle }O'A'B'&=\frac12|\vv{O'A'}\times\vv{O'B'}|\\
&=\frac5{18}\left|\cvv{1}{2}{2}\times\cvv{2}{1}{2}\right|\qquad(|\mathbf a\times(-\mathbf b)|=|\mathbf a\times\mathbf b|)\\
&=\frac5{18}\left|\cvv{2}{2}{-3}\right|\\
&=\frac5{18}\sqrt{17}\text{ units}^2.
\end{align*}

\clearpage
\par\noindent\begin{minipage}{\linewidth}
\Q{\textcolor{black!60}{\small[NYJC 2024 Promo Q5]}\par\smallskip
The position vectors of points $A$, $B$ and $C$, relative to the origin $O$ are $\mathbf a$, $\mathbf b$ and $\mathbf c$ respectively, where
\[\mathbf a=3\mathbf i+2\mathbf j-\mathbf k,\quad \mathbf b=4\mathbf i-3\mathbf j+2\mathbf k,\quad \mathbf c=3\mathbf i-\mathbf j-\mathbf k.\]
\begin{enumerate}[label=\textbf{(\alph*)},leftmargin=2.2em,itemsep=3pt]
\item Find $\mathbf a\times\mathbf b$ and deduce the exact area of triangle $OAB$.\hfill[3]
\item Hence, find the volume of the tetrahedron $OABC$.\par
[The volume of a tetrahedron is $\dfrac13\times\text{base area}\times\text{perpendicular height}$.]\hfill[3]
\item Find the perpendicular distance from $B$ to line $OA$.\hfill[2]
\end{enumerate}}
\end{minipage}\par
\par\medskip\hrule\medskip
\textbf{Solution}\par\smallskip
\textbf{(a)}
\begin{align*}
\mathbf a\times\mathbf b&=\cvv{3}{2}{-1}\times\cvv{4}{-3}{2}=\cvv{1}{-10}{-17}\\
\text{Area of }\triangle AOB&=\frac12|\mathbf a\times\mathbf b|\\
&=\frac12\left|\cvv{1}{-10}{-17}\right|\\
&=\frac12\sqrt{1+100+289}=\frac12\sqrt{390}\text{ units}^2.
\end{align*}
\textbf{(b)}
\begin{align*}
\text{Perpendicular height of tetrahedron}
&=\frac1{\sqrt{1+100+289}}\left|\vv{OC}\cdot\cvv{1}{-10}{-17}\right|\\
&=\frac1{\sqrt{390}}\left|\cvv{3}{-1}{-1}\cdot\cvv{1}{-10}{-17}\right|\\
&=\frac{30}{\sqrt{390}}\\
\text{Volume of tetrahedron }OABC
&=\frac13\left(\frac{\sqrt{390}}2\right)\left(\frac{30}{\sqrt{390}}\right)=5\text{ units}^3.
\end{align*}
\textbf{(c)}
\begin{align*}
\text{Perpendicular distance from }B\text{ to line }OA
&=|\mathbf b\times\hat{\mathbf a}|\\
&=\frac{\sqrt{390}}{\sqrt{9+4+1}}\\
&=\sqrt{\frac{195}{7}}\quad\text{or }5.28.
\end{align*}
\clearpage
\par\noindent\begin{minipage}{\linewidth}
\Q{\textcolor{black!60}{\small[MI 2024 Promo Q4]}\par\smallskip
Referred to the origin $O$, the points $A$, $B$ and $C$ have position vectors $\mathbf a$, $\mathbf b$ and $\mathbf c$. $C$ lies on $AB$ such that $AC:AB=1:3$.
\begin{enumerate}[label=\textbf{(\roman*)},leftmargin=2.2em,itemsep=3pt]
\item Find $\mathbf c$ in terms of $\mathbf a$ and $\mathbf b$.\hfill[2]
\item Show that the area of triangle $OAC$ is given by $k|\mathbf a\times\mathbf b|$, where $k$ is a constant to be determined.\hfill[2]
\end{enumerate}
It is now given that $\mathbf a$ and $\mathbf b$ are unit vectors.
\begin{enumerate}[label=\textbf{(\roman*)},leftmargin=2.2em,itemsep=3pt,start=3]
\item Using the result in part (i), find the value of $(2\mathbf a-\mathbf b)\cdot\mathbf c$.\hfill[3]
\end{enumerate}}
\end{minipage}\par
\par\medskip\hrule\medskip
\textbf{Solution}\par\smallskip
\textbf{(i)}
\begin{align*}
AC:AB&=1:3\quad\Rightarrow\quad AC:CB=1:2\\
\mathbf c&=\frac{(1)\mathbf b+(2)\mathbf a}{1+2}=\frac{2\mathbf a+\mathbf b}{3}.
\end{align*}
\textbf{(ii)}
\begin{align*}
\text{Area of triangle }OAC
&=\frac13\text{ Area of triangle }OAB\\
&=\frac13\times\frac12|\vv{OA}\times\vv{OB}|\\
&=\frac16|\mathbf a\times\mathbf b|,\qquad k=\frac16.
\end{align*}
\textbf{(iii)}
\begin{align*}
(2\mathbf a-\mathbf b)\cdot\mathbf c
&=(2\mathbf a-\mathbf b)\cdot\left(\frac{2\mathbf a+\mathbf b}{3}\right)\\
&=\frac13(2\mathbf a-\mathbf b)\cdot(2\mathbf a+\mathbf b)\\
&=\frac13(4\mathbf a\cdot\mathbf a+2\mathbf a\cdot\mathbf b-2\mathbf b\cdot\mathbf a-\mathbf b\cdot\mathbf b)\\
&=\frac13(4|\mathbf a|^2-|\mathbf b|^2)\\
&=\frac13\left(4(1)^2-(1)^2\right)\quad\text{since $\mathbf a$, $\mathbf b$ are unit vectors}\\
&=1.
\end{align*}
\clearpage
\par\noindent\begin{minipage}{\linewidth}
\Q{\textcolor{black!60}{\small[TMJC 2024 Promo Q6]}\par\smallskip
\begin{enumerate}[label=\textbf{(\alph*)},leftmargin=2.2em,itemsep=4pt]
\item The points $A$, $B$ and $C$ have position vectors $\cvv{1}{3}{-7}$, $\cvv{5}{1}{1}$ and $\cvv{-9}{8}{-27}$, with respect to the origin $O$.
\begin{enumerate}[label=(\roman*),leftmargin=1.8em,itemsep=3pt]
\item Show that the points $A$, $B$ and $C$ are collinear.\hfill[2]
\item Find the exact area of the triangle $OAB$.\hfill[3]
\end{enumerate}
\item The vectors $\mathbf u$, $\mathbf v$ and $\mathbf w$ are defined such that $|\mathbf u|=1$, $|\mathbf v|=2$ and $|\mathbf w|=3$.
\begin{enumerate}[label=(\roman*),leftmargin=1.8em,itemsep=3pt]
\item Suppose that $\mathbf u+\mathbf v=\mathbf w$. With the aid of a diagram, or otherwise, explain why $\mathbf u$ is a unit vector that is parallel to $\mathbf w$.\hfill[2]
\item Given instead that each of the three vectors is perpendicular to the sum of the other two vectors, find the exact value of $|\mathbf u+\mathbf v+\mathbf w|$.\hfill[4]
\end{enumerate}\end{enumerate}}

\end{minipage}\par
\par\medskip\hrule\medskip
\textbf{Solution}\par\smallskip
\textbf{(a)(i)}
\begin{align*}
\vv{AB}&=\cvv{5}{1}{1}-\cvv{1}{3}{-7}=\cvv{4}{-2}{8}\\
\vv{AC}&=\cvv{-9}{8}{-27}-\cvv{1}{3}{-7}=\cvv{-10}{5}{-20}
=-\frac52\cvv{4}{-2}{8}=-\frac52\vv{AB}.
\end{align*}
Since $\vv{AB}$ is parallel to $\vv{AC}$, with a common point $A$, the points $A$, $B$ and $C$ are collinear. (Shown)\par
\textbf{(a)(ii)}
\begin{align*}
\text{Area of triangle }OAB&=\frac12|\vv{OA}\times\vv{OB}|\\
&=\frac12\left|\cvv{1}{3}{-7}\times\cvv{5}{1}{1}\right|\\
&=\frac12\left|\cvv{10}{-36}{-14}\right|=\sqrt{398}.
\end{align*}
\textbf{(b)(i)} Since $\mathbf u+\mathbf v=\mathbf w$, and $|\mathbf u|=1$, $|\mathbf v|=2$ and $|\mathbf w|=3$, $\mathbf u$ and $\mathbf v$ must both be parallel to $\mathbf w$, as seen in the diagram.
\begin{center}\begin{tikzpicture}[>=Stealth,font=\small]
\draw[->,thick](0,0)--(1.5,0)node[midway,above]{$|\mathbf u|=1$};
\draw[->,thick](1.5,0)--(4.5,0)node[midway,above]{$|\mathbf v|=2$};
\draw[->,thick](0,-.5)--(4.5,-.5)node[midway,below]{$|\mathbf w|=3$};
\end{tikzpicture}\end{center}
Since $|\mathbf u|=1$, $\mathbf u$ is a unit vector parallel to $\mathbf w$.\par
\textbf{(b)(ii)} Since each of the three vectors is perpendicular to the sum of the other two vectors,
\begin{align*}
\mathbf u\cdot(\mathbf v+\mathbf w)&=0,\qquad\mathbf v\cdot(\mathbf u+\mathbf w)=0,\qquad\mathbf w\cdot(\mathbf v+\mathbf u)=0\\
|\mathbf u+\mathbf v+\mathbf w|^2
&=(\mathbf u+\mathbf v+\mathbf w)\cdot(\mathbf u+\mathbf v+\mathbf w)\\
&=\mathbf u\cdot\mathbf u+\mathbf v\cdot\mathbf v+\mathbf w\cdot\mathbf w+\mathbf u\cdot(\mathbf v+\mathbf w)+\mathbf v\cdot(\mathbf u+\mathbf w)+\mathbf w\cdot(\mathbf v+\mathbf u)\\
&=|\mathbf u|^2+|\mathbf v|^2+|\mathbf w|^2\\
&=1^2+2^2+3^2=14.
\end{align*}
Since $|\mathbf u+\mathbf v+\mathbf w|>0$, $|\mathbf u+\mathbf v+\mathbf w|=\sqrt{14}$.

\clearpage
\par\noindent\begin{minipage}{\linewidth}
\Q{\textcolor{black!60}{\small[RI 2025 Prelim P1 Q3]}\par\smallskip
Relative to the origin $O$, the points $A$ and $C$ have position vectors $\mathbf a$ and $2\mathbf b$ such that
\[\mathbf a=5p\mathbf i-2p\mathbf j+4p\mathbf k\quad\text{and}\quad\mathbf b=\mathbf i-2\mathbf j+2\mathbf k,\]
where $p$ is a positive constant. It is given that $|\mathbf a|=2|\mathbf b|$.
\begin{enumerate}[label=\textbf{(\alph*)},leftmargin=2.2em,itemsep=3pt]
\item Find the exact value of $p$.\hfill[2]
\item Evaluate $(\mathbf a+2\mathbf b)\cdot(\mathbf a-2\mathbf b)$.\hfill[1]
\item Evaluate $\mathbf a\times\mathbf b$ and hence find the area of triangle $OAC$.\hfill[3]
\item Use a geometrical reason to explain why $|\mathbf a\times\mathbf b|=\dfrac14|\mathbf a+2\mathbf b||\mathbf a-2\mathbf b|$.\hfill[2]
\end{enumerate}}

\end{minipage}\par
\par\medskip\hrule\medskip
\textbf{Solution}\par\smallskip
\textbf{(a)} Since $|\mathbf a|=2|\mathbf b|$,
\begin{align*}
\sqrt{(5p)^2+(-2p)^2+(4p)^2}&=2\sqrt{1^2+(-2)^2+2^2}\\
3\sqrt5p&=6\qquad(p>0)\\
p&=\frac2{\sqrt5}=\frac{2\sqrt5}{5}.
\end{align*}
\textbf{(b)}
\begin{align*}
(\mathbf a+2\mathbf b)\cdot(\mathbf a-2\mathbf b)
&=\mathbf a\cdot\mathbf a-2\mathbf a\cdot\mathbf b+2\mathbf b\cdot\mathbf a-4\mathbf b\cdot\mathbf b\\
&=|\mathbf a|^2-4|\mathbf b|^2\qquad(|\mathbf a|=2|\mathbf b|)\\
&=0.
\end{align*}
\textbf{(c)}
\begin{align*}
\mathbf a\times\mathbf b&=\cvv{5p}{-2p}{4p}\times\cvv{1}{-2}{2}
=p\cvv{5}{-2}{4}\times\cvv{1}{-2}{2}
=p\cvv{4}{-6}{-8}=\frac{4\sqrt5}{5}\cvv{2}{-3}{-4}\\
\text{Area of triangle }OAC&=\frac12|\mathbf a\times2\mathbf b|=|\mathbf a\times\mathbf b|\\
&=\frac{4\sqrt5}{5}\sqrt{2^2+(-3)^2+(-4)^2}=\frac{4\sqrt{145}}5.
\end{align*}
\textbf{(d)} Observe that $\vv{OD}=\mathbf a+2\mathbf b$ and $\vv{CA}=\mathbf a-2\mathbf b$. From (b), $OD$ is perpendicular to $CA$. In addition, $CA$ and $OD$, which are the diagonals of the parallelogram $OADC$, bisect each other.
\begin{center}\begin{tikzpicture}[scale=.8,font=\small,>=Stealth]
\coordinate(O)at(0,0);\coordinate(A)at(4,0);\coordinate(C)at(2,3.464);\coordinate(D)at(6,3.464);
\draw[semithick](O)--(A)--(D)--(C)--cycle;
\node[below left]at(O){$O$};\node[below right]at(A){$A$};\node[above left]at(C){$C$};\node[above right]at(D){$D$};
\draw[->](O)--(2,0)node[midway,below]{$\mathbf a$};\draw[->](O)--(1,1.732)node[midway,left]{$2\mathbf b$};
\end{tikzpicture}\end{center}
From (c), area of triangle $OAC=|\mathbf a\times\mathbf b|$. But area of triangle $OAC$ is also
\begin{align*}
\frac12(CA)\left(\frac12OD\right)&=\frac14|\mathbf a-2\mathbf b||\mathbf a+2\mathbf b|.
\end{align*}
Hence $|\mathbf a\times\mathbf b|=\dfrac14|\mathbf a+2\mathbf b||\mathbf a-2\mathbf b|$.

\par\vspace{\baselineskip}\begin{center}\textbf{END OF TEST}\end{center}\end{document}
